% !PRINTED 599
\weberpage{599}
%<<CONT
whether one or several of the roots can be represented in this way,
while others admit no such representation. This question is justified
only for irreducible equations, since for reducible equations all
conceivable combinations can occur; and here the following theorem now holds:

\begin{quote}
\noindent IV.\ \gesperrt{If a root of an irreducible equation is
determinable by the solution of cyclic equations, then the equation is metacyclic.}
\end{quote}

If an irreducible equation of degree $n$, $f(x) = 0$, has even
one root which becomes rational through successive adjunction of roots
of cyclic equations, then it must necessarily become reducible
through this adjunction, since in the end a linear
factor must separate off. Let $P$ then be the group
of our equation, after everything necessary has been adjoined so far
that $f(x)$ is not yet reducible, but splits into factors through the next
adjunction of the root $\varepsilon$ of a cyclic equation of prime
degree $m$. Then by \S.\,175, $P$ must have
a normal divisor $Q$ of index $m$, to which the
group of the equation reduces after the adjunction of $\varepsilon$, and the
permutation group $Q$ must be intransitive. If $Q$ is the unit
group, then $f(x) = 0$ is completely solved by the adjunction of $\varepsilon$. But if $Q$ is still different from the unit group, then the theorems of \S.\,158 come into force. According to these,
$f(x)$ splits after the adjunction of $\varepsilon$ — since the index of $Q$ is here a
prime — into $m$ factors, $\varphi(x, \varepsilon), \varphi_1(x, \varepsilon), \ldots \varphi_{m-1}(x, \varepsilon)$, of
the same degree $\mu$, and we have $\mu m = n$.

If, for example, $n$ is a prime, then $\mu = 1$ must hold; the functions $\varphi$ are linear, and the equation $f(x) = 0$ is completely solved; hence Theorem IV is correct for equations of prime degree. In the general case, one of the factors
of degree $\mu$, $\varphi(x, \varepsilon), \varphi_1(x, \varepsilon), \ldots \varphi_{m-1}(x, \varepsilon)$, say $\varphi(x, \varepsilon)$, must have a
root which becomes rational through the adjunction of the roots of cyclic
equations; and if we therefore assume that
our theorem is already proved for equations of degree $\mu$, it
follows that $\varphi(x, \varepsilon) = 0$ itself, and hence also its group, is metacyclic.

But since by \S.\,158 the different equations $\varphi = 0$,
$\varphi_1 = 0, \ldots \varphi_{m-1} = 0$ have the same group, they are all,
and consequently also $f(x) = 0$, metacyclic.
